%% ============================================================
%%  Aula 5 — Gradientes de política I
%%  Adaptação em pt-BR do CS234 (Stanford, Emma Brunskill), Lecture 5
%%  Prof. Marcelo Keese Albertini — Universidade Federal de Uberlândia
%%  Compilar com XeLaTeX (2x) — latexmk -xelatex aula05.tex
%% ============================================================
\documentclass[aspectratio=169, 11pt]{beamer}
%% .sty do projeto na raiz do repositório (../)
\makeatletter
\def\input@path{{./}{../}}
\makeatother


\usetheme{AprendizadoRL}      % ANTES do babel — fontspec precisa vir primeiro
\usepackage{babel}
\babelprovide[import, main]{brazilian}
\usepackage{booktabs}

\title[Aula 5 — Gradientes de política I]{Aula 5: Gradientes de política I}
\subtitle[Aprendizagem por Reforço]{Políticas parametrizadas \textbullet\ Razão de verossimilhança \textbullet\ REINFORCE \textbullet\ Linha de base e vantagem}
\author[Prof. Marcelo K. Albertini]{Prof. Marcelo Keese Albertini \\
  \footnotesize Universidade Federal de Uberlândia \\
  \footnotesize adaptado de Emma Brunskill, CS234 (Stanford)}
\institute{Universidade Federal de Uberlândia}
\date{2026}

% ---- Nota discreta: perguntas ficam para o professor conduzir em aula
\newcommand{\paradiscussao}{\smallskip
  {\footnotesize\color{rlCinza}\itshape Pergunta para discussão conduzida em aula.}}

%% ------------------------------------------------------------
%% Macros locais desta aula
%% ------------------------------------------------------------
\newcommand{\thv}{\theta}                       % vetor de parâmetros da política
\newcommand{\wv}{w}                             % pesos do crítico / valor
\newcommand{\polth}{\pol_{\thv}}                % política parametrizada
\newcommand{\gth}{\nabla_{\thv}}                % gradiente em theta
\newcommand{\traj}{\tau}                        % trajetória
\newcommand{\Rtraj}{R(\traj)}
\newcommand{\Ptraj}{P(\traj;\thv)}
\newcommand{\Vhat}{\hat{V}}
\newcommand{\Qhat}{\hat{Q}}
\newcommand{\Ahat}{\hat{A}}
\newcommand{\Ghat}{\hat{G}}
\newcommand{\Adv}{A}                            % função vantagem
\newcommand{\Advpi}{A^{\pol}}
\newcommand{\feat}{\phi}                        % vetor de atributos
\newcommand{\gscore}{\gth \log \polth(a \mid s)} % função escore
\newcommand{\Var}{\operatorname{Var}}

\begin{document}

\begin{frame}[plain, noframenumbering]\titlepage\end{frame}

%% ------------------------------------------------------------
\begin{frame}{Onde estamos}{Do valor à política}

  \begin{itemize}\setlength{\itemsep}{0.5ex}
    \item \textbf{Aulas 2--3:} com um modelo do mundo, calcular $\Vstar$ e
      $\polstar$ por programação dinâmica.
    \item \textbf{Aula 4:} sem modelo, \emph{aprender} $\Qhat(s,a;\wv)$ a partir
      de amostras --- MC, SARSA, Q-learning, aproximação de função, DQN.
    \item \textbf{Hoje:} descartar o intermediário. Parametrizar
      \emph{diretamente} a política $\polth(a\mid s)$ e subir o gradiente do
      retorno esperado.
    \item \textbf{Próxima aula:} tornar esse passo de gradiente confiável ---
      gradiente natural, TRPO, PPO.
  \end{itemize}

  \begin{ideiabox}
    Em todo o curso até aqui, a política foi um \emph{subproduto} do valor:
    $\pol(s) = \argmax_a \Qhat(s,a)$. Hoje ela vira o objeto de primeira
    classe, e o valor --- quando aparecer --- será apenas um auxiliar para
    reduzir variância.
  \end{ideiabox}
\end{frame}

%% ------------------------------------------------------------
\begin{frame}{Recapitulação: o que ficou da aula passada}

  Todo método de controle sem modelo da Aula 4 tinha a mesma forma ---
  atualização incremental rumo a um \emph{alvo}:
  \[
    \Delta \wv = \alpha\bigl(\text{alvo} - \Qhat(s,a;\wv)\bigr)\, \nabla_{\wv} \Qhat(s,a;\wv)
  \]

  \begin{center}\small
  \begin{tabular}{@{}ll@{}}
    \toprule
    \textbf{Método} & \textbf{Alvo} \\
    \midrule
    Monte Carlo   & $G_t$ \\
    SARSA         & $r + \gamma\,\Qhat(s',a';\wv)$ \\
    Q-learning    & $r + \gamma \max_{a'} \Qhat(s',a';\wv)$ \\
    DQN           & $r + \gamma \max_{a'} \Qhat(s',a';\wv^{-})$, amostrado de um \emph{buffer} \\
    \bottomrule
  \end{tabular}
  \end{center}

  \begin{alertabox}
    E a \textbf{tríade mortal} --- aproximação de função $+$
    \emph{bootstrapping} $+$ off-policy --- que retira as garantias. O DQN não
    a resolve: contorna-a com repetição de experiência e alvos fixos.
  \end{alertabox}
\end{frame}

%% ------------------------------------------------------------
\section{Por que parametrizar a política}

%% ------------------------------------------------------------
\begin{frame}{Três famílias de algoritmos}

  \begin{center}
  \scalebox{0.95}{\begin{tikzpicture}[font=\small, node distance=6mm]
    \node[draw=rlAzul, very thick, fill=rlAzul!8, rounded corners=3pt,
          text width=33mm, align=center, minimum height=20mm] (v)
      {\textbf{Baseado em valor}\\[0.3em]
       {\scriptsize aprende $\Qhat$\\ política \emph{implícita}\\ ($\epsilon$-guloso)}};
    \node[draw=rlLaranja, very thick, fill=rlLaranja!10, rounded corners=3pt,
          text width=33mm, align=center, minimum height=20mm, right=26mm of v] (p)
      {\textbf{Baseado em política}\\[0.3em]
       {\scriptsize aprende $\polth$\\ \emph{sem} função valor}};
    \node[draw=rlTeal, very thick, fill=rlTeal!10, rounded corners=3pt,
          text width=33mm, align=center, minimum height=13mm,
          below=9mm of $(v)!0.5!(p)$] (ac)
      {\textbf{Ator--crítico}\\[0.2em]
       {\scriptsize aprende \emph{ambos}}};
    \draw[trans] (v.south) -- (ac.west);
    \draw[trans] (p.south) -- (ac.east);
  \end{tikzpicture}}
  \end{center}

  \begin{itemize}\setlength{\itemsep}{0.3ex}
    \item Hoje construímos a coluna do meio; o ator--crítico aparece no fim da
      aula, quase de graça.
    \item Notação: $\wv$ continua sendo os pesos de uma \emph{função valor};
      $\thv$ passa a ser os pesos da \emph{política}.
  \end{itemize}
\end{frame}

%% ------------------------------------------------------------
\begin{frame}{A política parametrizada}

  \begin{defbox}{Política estocástica parametrizada}
    \[
      \polth(a \mid s) = \Prob[a_t = a \mid s_t = s;\, \thv],
      \qquad \thv \in \R^{n}
    \]
    Uma distribuição sobre ações, diferenciável em $\thv$ sempre que é não
    nula. O objetivo é achar $\thv$ que maximize o valor da política.
  \end{defbox}

  \begin{itemize}\setlength{\itemsep}{0.4ex}
    \item Em espaços discretos: uma rede que produz \emph{logits} e uma
      \emph{softmax} no fim.
    \item Em espaços contínuos: uma rede que produz a média (e talvez o
      desvio) de uma gaussiana.
    \item Continuamos \textbf{sem modelo}: nada de $\tran{s'}{s,a}$ será
      necessário --- e isso vai cair sozinho da matemática, o que é um dos
      fatos mais bonitos da aula.
  \end{itemize}
\end{frame}

%% ------------------------------------------------------------
\begin{frame}{Por que não bastam políticas determinísticas}{Três motivos}

  \begin{enumerate}\setlength{\itemsep}{0.55ex}
    \item \textbf{Ações contínuas.} O $\argmax_a \Qhat(s,a)$ vira um problema
      de otimização \emph{dentro} de cada passo. Com $\polth$ a ação sai por
      uma passagem para frente na rede.
    \item \textbf{Estados perceptualmente ambíguos.} Se o agente não
      distingue dois estados, a melhor política pode exigir
      \emph{aleatoriedade} --- veremos o exemplo já a seguir.
    \item \textbf{Suavidade.} Um $\epsilon$-guloso muda de ação de forma
      descontínua quando dois valores $\Qhat$ se cruzam; $\polth$ muda
      suavemente com $\thv$. Isso costuma dar convergência mais estável.
  \end{enumerate}

  \begin{ideiabox}
    Há um quarto motivo, moderno: quando a política é um modelo de linguagem,
    ela \emph{já é} uma distribuição sobre tokens. Ajustá-la por gradiente de
    política é a coisa mais natural do mundo --- é o que fazem RLHF, PPO e
    GRPO.
  \end{ideiabox}
\end{frame}

%% ------------------------------------------------------------
\begin{frame}{Exemplo: o corredor com estados sinônimos}

  Corredor de cinco casas. O agente enxerga apenas atributos locais do tipo
  \[
    \feat(s,a) = \mathbf{1}\{\text{há parede ao N},\; a = \text{mover para L}\}
    \quad \text{(e análogos para N, S, L, O)}
  \]

  \begin{center}
  \scalebox{0.9}{\begin{tikzpicture}[font=\small]
    \foreach \i/\fill/\lab in {1/white/, 2/rlCinza!35/, 3/white/, 4/rlCinza!35/, 5/white/}{
      \node[draw=rlAzul, thick, fill=\fill, minimum size=13mm]
        (c\i) at (\i*14mm, 0) {};}
    \node at (c3) {\textcolor{rlVerde!70!black}{\Large\textbf{\$}}};
    \node at (c1) {\textcolor{rlVermelho}{\Large$\times$}};
    \node at (c5) {\textcolor{rlVermelho}{\Large$\times$}};
    \node[font=\scriptsize, below=1mm of c2, text=rlCinza] {idênticos};
    \node[font=\scriptsize, below=1mm of c4, text=rlCinza] {idênticos};
    \node[font=\scriptsize, above=1mm of c3, text=rlVerde!70!black] {objetivo};
    \draw[rlAzul, very thick] ($(c1.north west)+(0,0.6mm)$) -- ($(c5.north east)+(0,0.6mm)$);
    \draw[rlAzul, very thick] ($(c1.south west)-(0,0.6mm)$) -- ($(c5.south east)-(0,0.6mm)$);
  \end{tikzpicture}}
  \end{center}

  \begin{itemize}\setlength{\itemsep}{0.3ex}
    \item As casas 2 e 4 têm paredes ao N e ao S: geram \textbf{exatamente o
      mesmo vetor de atributos}. O agente não pode distingui-las.
    \item Comparemos duas abordagens: valor aproximado
      $Q_{\thv}(s,a) = f(\feat(s,a);\thv)$ contra política parametrizada
      $\polth(a\mid s) = g(\feat(s,a);\thv)$.
  \end{itemize}
\end{frame}

%% ------------------------------------------------------------
\begin{frame}{O corredor: por que o determinismo falha}

  Como 2 e 4 são indistinguíveis, \emph{qualquer} política determinística é
  obrigada a escolher a mesma ação nas duas:

  \begin{columns}[T]
  \begin{column}{0.5\textwidth}
    \begin{alertabox}
      \textbf{Sempre para o Leste:} da casa 2 chega ao objetivo; da casa 4 vai
      para a 5 e fica preso. \textbf{Sempre para o Oeste:} simétrico. Em ambos
      os casos, \emph{metade dos inícios nunca alcança o objetivo}.
    \end{alertabox}
  \end{column}
  \begin{column}{0.5\textwidth}
    \begin{ideiabox}
      RL baseado em valor aprende uma política quase determinística. O
      $\epsilon$ salva no limite, mas o tempo esperado para escapar cresce
      como $1/\epsilon$: milhares de passos.
    \end{ideiabox}
  \end{column}
  \end{columns}

  \begin{teobox}{A política ótima aqui é estocástica}
    $\polth(\text{L}\mid \text{parede N e S}) = \polth(\text{O}\mid \text{parede N e S}) = 0{,}5$:
    um passeio aleatório local, que alcança o objetivo em poucos passos com
    alta probabilidade a partir de qualquer início.
  \end{teobox}
\end{frame}

%% ------------------------------------------------------------
\begin{frame}{Teste sua compreensão}

  \begin{quizbox}{}
    O exemplo do corredor mostra que políticas estocásticas podem ser
    estritamente melhores que determinísticas.

    \smallskip
    Isso contradiz o teorema da Aula 2, que garante existir uma política
    \emph{determinística} ótima para todo MDP finito?
    \paradiscussao
  \end{quizbox}

  \pause

  \begin{respbox}
    \textbf{Não há contradição --- o objeto mudou.} O teorema vale para MDPs,
    em que a política é função do \emph{estado}. Aqui a política é função da
    \emph{observação} $\feat(s,\cdot)$, que não identifica o estado: estamos
    de fato num POMDP (ou, equivalentemente, numa classe restrita de políticas
    que não separa 2 de 4). Sob restrição de representação --- ou sob
    observabilidade parcial --- o ótimo pode exigir aleatoriedade.
  \end{respbox}
\end{frame}

%% ------------------------------------------------------------
\section{O objetivo e o truque da razão de verossimilhança}

%% ------------------------------------------------------------
\begin{frame}{O que exatamente vamos maximizar}

  \begin{defbox}{Valor de uma política parametrizada}
    Para MDPs episódicos com estado inicial $s_0$ e horizonte $T$,
    \[
      V(\thv) \defeq V(s_0,\thv)
      = \E_{\polth}\Bigl[\; \sum_{t=0}^{T} \rec(s_t,a_t) \;\Bigm|\; \polth,\, s_0 \Bigr]
    \]
    \vspace{-0.6em}
    A esperança é sobre estados \emph{e} ações visitados por $\polth$.
  \end{defbox}

  Duas reescritas equivalentes --- a escolha importa:
  \begin{itemize}\setlength{\itemsep}{0.2ex}
    \item \textbf{Por ação:} $V(s_0,\thv) = \sum_a \polth(a\mid s_0) Q(s_0,a,\thv)$
      \hfill{\scriptsize (Sutton \& Barto, 13.2)}
    \item \textbf{Por trajetória:} $V(\thv) = \sum_{\traj} \Ptraj\, \Rtraj$
      \hfill{\scriptsize (o caminho de hoje)}
  \end{itemize}
  \vspace{-0.3em}

  \vspace{-0.6em}
  \begin{ideiabox}
    A segunda forma trata tudo como \emph{otimização estocástica pura}:
    derivação curta, e válida bem além de MDPs.
  \end{ideiabox}
\end{frame}

%% ------------------------------------------------------------
\begin{frame}{Notação de trajetória}

  \begin{defbox}{Trajetória, retorno e verossimilhança}
    \[
      \traj = (s_0,a_0,r_0,\; s_1,a_1,r_1,\; \dots,\; s_{T-1},a_{T-1},r_{T-1},s_T)
    \]
    \[
      \Rtraj = \sum_{t=0}^{T} \rec(s_t,a_t),
      \qquad
      \Ptraj = \mu(s_0) \prod_{t=0}^{T-1}
        \underbrace{\polth(a_t\mid s_t)}_{\text{política}}
        \underbrace{\tran{s_{t+1}}{s_t,a_t}}_{\text{dinâmica}}
    \]
  \end{defbox}

  O problema de otimização fica
  $\displaystyle \argmax_{\thv} V(\thv) = \argmax_{\thv} \sum_{\traj} \Ptraj\, \Rtraj$.
  \smallskip

  \begin{alertabox}
    Note quem depende de $\thv$: \textbf{apenas} os fatores
    $\polth(a_t\mid s_t)$ --- nem $\mu$, nem a dinâmica, nem $\rec$.
  \end{alertabox}
\end{frame}

%% ------------------------------------------------------------
\begin{frame}{Antes do gradiente: por que não otimização sem gradiente?}

  Maximizar $V(\thv)$ é otimização de caixa-preta. Existem métodos que
  \emph{não} usam gradiente:

  \begin{itemize}\setlength{\itemsep}{0.2ex}
    \item \textbf{Busca aleatória}: perturbar $\thv$, manter se melhorou.
    \item \textbf{Diferenças finitas}:
      $\partial V/\partial \thv_k \approx [V(\thv + \varepsilon e_k) - V(\thv)]/\varepsilon$.
    \item \textbf{Estratégias evolutivas} e \emph{cross-entropy method}.
  \end{itemize}

  \begin{teobox}{Por que insistir no gradiente}
    Diferenças finitas exigem $\mathcal{O}(n)$ avaliações \emph{ruidosas} por
    passo --- proibitivo com $n$ na casa dos milhões. O gradiente analítico
    custa \emph{uma} passagem para trás, qualquer que seja $n$.
  \end{teobox}

  \begin{ideiabox}
    Não os descarte: são competitivos quando $n$ é pequeno ou quando o
    simulador é barato e massivamente paralelo (Salimans et al., 2017).
  \end{ideiabox}
\end{frame}

%% ------------------------------------------------------------
\begin{frame}{Subida de gradiente na política}

  \begin{defbox}{Gradiente de política}
    \[
      \Delta \thv = \alpha\, \gth V(s_0,\thv),
      \qquad
      \gth V(s_0,\thv) =
      \begin{pmatrix}
        \partial V(s_0,\thv) / \partial \thv_1 \\[-0.2em]
        \vdots \\[-0.2em]
        \partial V(s_0,\thv) / \partial \thv_n
      \end{pmatrix}
    \]
    Sobe-se esse gradiente rumo a um \textbf{máximo local}; $\alpha$ é o passo.
  \end{defbox}

  \begin{alertabox}
    Local, não global --- a principal desvantagem da família. Em troca:
    convergência suave, ações contínuas e políticas estocásticas.
  \end{alertabox}

  A dificuldade: $V(\thv) = \sum_{\traj}\Ptraj \Rtraj$ soma sobre um número
  astronômico de trajetórias, e não sabemos calcular $\Ptraj$.
\end{frame}

%% ------------------------------------------------------------
\begin{frame}{O truque da razão de verossimilhança}

  {\small
  \begin{align*}
    \gth V(\thv)
      &= \gth \sum_{\traj} \Ptraj\, \Rtraj
       \;=\; \sum_{\traj} \gth \Ptraj\; \Rtraj
        &&\text{\scriptsize ($\Rtraj$ não depende de $\thv$)} \\[-0.1em]
    \onslide<2->{
      &= \sum_{\traj} \frac{\Ptraj}{\Ptraj}\, \gth \Ptraj\; \Rtraj
        &&\text{\scriptsize (multiplicar e dividir por $\Ptraj$)} \\[-0.1em]}
    \onslide<3->{
      &= \sum_{\traj} \Ptraj\, \Rtraj\,
         \underbrace{\frac{\gth \Ptraj}{\Ptraj}}_{\text{razão de verossimilhança}}
       \;=\; \sum_{\traj} \Ptraj\, \Rtraj\, \gth \log \Ptraj
        &&\text{\scriptsize (pois $\nabla \log f = \nabla f / f$)}}
  \end{align*}}

  \onslide<3->{
  \begin{teobox}{O gradiente virou uma esperança}
    $\gth V(\thv) = \E_{\traj \sim \polth}\bigl[\Rtraj\; \gth \log \Ptraj\bigr]$
    --- e esperanças a gente estima por amostragem.
  \end{teobox}}
\end{frame}

%% ------------------------------------------------------------
\begin{frame}{Estimador de Monte Carlo}

  Com $m$ trajetórias $\traj^{(1)},\dots,\traj^{(m)}$ coletadas rodando a
  política atual:
  \[
    \gth V(\thv) \;\approx\; \hat{g}
      = \frac{1}{m}\sum_{i=1}^{m} R(\traj^{(i)})\, \gth \log P(\traj^{(i)};\thv)
  \]

  \begin{itemize}\setlength{\itemsep}{0.4ex}
    \item O estimador é \textbf{não enviesado}: $\E[\hat{g}] = \gth V(\thv)$
      exatamente, para qualquer $m$.
    \item Mas ainda parece inútil: $\log P(\traj;\thv)$ contém a dinâmica, que
      é desconhecida.
    \item \textbf{Próximo slide:} ao abrir o logaritmo, a dinâmica desaparece.
  \end{itemize}

  \begin{ideiabox}
    Repare no que já conseguimos: transformamos ``derivar uma soma sobre todas
    as trajetórias'' em ``rodar a política algumas vezes e tirar uma média''.
    Todo o resto da aula é reduzir a \emph{variância} dessa média.
  \end{ideiabox}
\end{frame}

%% ------------------------------------------------------------
\begin{frame}{Decompondo a trajetória: a dinâmica cai fora}

  \vspace{-0.4em}
  {\small
  \begin{align*}
    \gth \log P(\traj^{(i)};\thv)
      &= \gth \log \Bigl[\,
         \underbrace{\mu(s_0)}_{\text{distrib.\ inicial}}
         \prod_{t=0}^{T-1}
         \underbrace{\polth(a_t\mid s_t)}_{\text{política}}\,
         \underbrace{\tran{s_{t+1}}{s_t,a_t}}_{\text{modelo de dinâmica}}
         \Bigr] \\[0.2em]
    \onslide<2->{
      &= \gth \Bigl[\, \log \mu(s_0)
         + \sum_{t=0}^{T-1} \log \polth(a_t\mid s_t)
         + \sum_{t=0}^{T-1} \log \tran{s_{t+1}}{s_t,a_t} \Bigr] \\[-0.1em]}
    \onslide<3->{
      &= \sum_{t=0}^{T-1}
         \underbrace{\gth \log \polth(a_t\mid s_t)}_{\text{função escore}}}
  \end{align*}}

  \onslide<3->{
  \begin{teobox}{Nenhum modelo de dinâmica é necessário}
    O logaritmo transforma produto em soma; a derivada aniquila o que não
    depende de $\thv$. Sobra só a política --- que \emph{nós} escolhemos.
  \end{teobox}}
\end{frame}

%% ------------------------------------------------------------
\begin{frame}{Gradiente de política por função escore}

  \begin{teobox}{Estimador completo (REINFORCE ingênuo)}
    \[
      \gth V(\thv) \;\approx\; \hat{g}
        = \frac{1}{m}\sum_{i=1}^{m} R(\traj^{(i)})
          \sum_{t=0}^{T-1} \gth \log \polth\bigl(a_t^{(i)} \bigm| s_t^{(i)}\bigr)
    \]
  \end{teobox}

  Tudo que é preciso para calculá-lo:
  \begin{enumerate}\setlength{\itemsep}{0.25ex}
    \item Rodar a política e guardar $(s_t,a_t,r_t)$.
    \item Somar as recompensas de cada trajetória.
    \item Derivar $\log \polth$ --- diferenciação automática faz isso.
  \end{enumerate}

  \begin{ideiabox}
    Sem modelo, sem $\max$ sobre ações. Em compensação: precisamos de
    \emph{episódios completos}, e o estimador é on-policy.
  \end{ideiabox}
\end{frame}

%% ------------------------------------------------------------
\begin{frame}{A função escore}

  \begin{defbox}{Função escore \textnormal{(\emph{score function})}}
    A derivada do logaritmo de uma verossimilhança parametrizada,
    $\gth \log p(x;\thv)$. Aqui, $\gscore$: a direção em $\thv$ que
    \emph{mais aumenta} a log-probabilidade de ter escolhido $a$ em $s$.
  \end{defbox}

  \begin{teobox}{Propriedade fundamental: o escore tem média zero}
    \vspace{-0.9em}
    \[
      \E_{a\sim\polth(\cdot\mid s)}\bigl[\gscore\bigr]
      = \sum_a \polth(a\mid s)\frac{\gth\polth(a\mid s)}{\polth(a\mid s)}
      = \gth \sum_a \polth(a\mid s) = \gth 1 = 0
    \]
    \vspace{-0.9em}

    Guarde esta identidade: é ela que provará, adiante, que uma
    \emph{linha de base} não introduz viés.
  \end{teobox}
\end{frame}

%% ------------------------------------------------------------
\begin{frame}{Política softmax e seu escore}

  Ações discretas, atributos $\feat(s,a)$, pesos lineares:
  \[
    \polth(a \mid s) = \frac{e^{\feat(s,a)^{\!\top}\thv}}
                            {\sum_{a'} e^{\feat(s,a')^{\!\top}\thv}}
  \]

  {\small
  \begin{align*}
    \gscore
      &= \gth\Bigl[\feat(s,a)^{\!\top}\thv - \log \sum_{a'} e^{\feat(s,a')^{\!\top}\thv}\Bigr] \\
      &= \feat(s,a) - \frac{\sum_{a'} \feat(s,a')\, e^{\feat(s,a')^{\!\top}\thv}}
                            {\sum_{a'} e^{\feat(s,a')^{\!\top}\thv}}
       \;=\; \feat(s,a) - \E_{a'\sim\polth}[\feat(s,a')]
  \end{align*}}

  \begin{ideiabox}
    \textbf{Atributo observado menos atributo esperado.} Se a ação tomada era
    ``mais típica'' do que a média, o escore é pequeno; se era atípica, o
    escore é grande. Multiplicado por um retorno alto, isso empurra a política
    justamente na direção do que foi surpreendentemente bom.
  \end{ideiabox}
\end{frame}

%% ------------------------------------------------------------
\begin{frame}{Política gaussiana e seu escore}

  Para ações contínuas, $a \in \R$ (ou $\R^d$):
  \[
    \mu_{\thv}(s) = \feat(s)^{\!\top}\thv,
    \qquad
    a \sim \mathcal{N}\bigl(\mu_{\thv}(s),\, \sigma^2\bigr)
  \]
  \[
    \gscore = \frac{\bigl(a - \mu_{\thv}(s)\bigr)\,\feat(s)}{\sigma^2}
  \]

  \begin{itemize}\setlength{\itemsep}{0.35ex}
    \item $\sigma^2$ pode ser fixa ou parametrizada --- na prática funciona
      como \textbf{controle de exploração}, que a otimização vai encolhendo.
    \item Redes profundas substituem a forma linear sem alterar a derivação.
  \end{itemize}

  \begin{alertabox}
    O fator $1/\sigma^2$ explode quando $\sigma \to 0$: uma política que ficou
    quase determinística produz gradientes gigantescos e instáveis. É um dos
    modos de falha mais comuns em controle contínuo --- limite $\sigma$ por
    baixo.
  \end{alertabox}
\end{frame}

%% ------------------------------------------------------------
\begin{frame}{Intuição do estimador por função escore}

  Forma genérica de cada termo, $\hat{g}_i = f(x_i)\, \gth \log p(x_i;\thv)$:

  \begin{center}
  \scalebox{0.74}{\begin{tikzpicture}[font=\scriptsize]
    \draw[rlCinza, thick, ->] (-0.3,0) -- (8.6,0) node[right] {$x$};
    \draw[domain=0:8, smooth, variable=\x, rlAzul, very thick]
      plot ({\x}, {2.1*exp(-(\x-3.4)*(\x-3.4)/2.2)});
    \node[rlAzul] at (1.3,1.55) {$p(x;\thv)$};
    \foreach \x/\f/\c in {1.9/0.35/rlCinza, 2.9/0.6/rlCinza, 3.6/1.5/rlLaranja,
                          4.4/1.05/rlAmbar, 5.4/0.3/rlCinza}{
      \draw[\c, line width=1.6pt, -{Stealth[length=2.4mm]}] (\x,-0.15) -- (\x,-0.15-\f*0.85);
      \fill[\c] (\x,{2.1*exp(-(\x-3.4)*(\x-3.4)/2.2)}) circle (1.6pt);}
    \node[rlLaranja, below] at (3.6,-1.5) {$f$ alto};
    \node[rlCinza, below] at (1.9,-0.55) {$f$ baixo};
    \draw[rlLaranja, line width=2pt, -{Stealth[length=3mm]}] (3.4,2.5) -- (4.9,2.5);
    \node[rlLaranja, right] at (5.0,2.5) {a densidade se desloca para cá};
  \end{tikzpicture}}
  \end{center}

  \begin{itemize}\setlength{\itemsep}{0.3ex}
    \item Andar na direção $\hat{g}_i$ \textbf{aumenta a log-probabilidade da
      amostra}, na proporção de quão boa ela foi.
    \item Vale mesmo que $f$ seja \textbf{descontínua, desconhecida ou
      discreta}: nunca derivamos $f$, só $\log p$.
  \end{itemize}
\end{frame}

%% ------------------------------------------------------------
\begin{frame}{Por que isso é notável: escore \emph{versus} reparametrização}

  \begin{columns}[T]
  \begin{column}{0.5\textwidth}
    \begin{defbox}{Escore (o que fizemos)}
      \vspace{-0.8em}
      \[ \gth \E_{x\sim p_{\thv}}[f(x)] = \E\bigl[f(x)\gth \log p_{\thv}(x)\bigr] \]
      \vspace{-1.6em}
      \begin{itemize}\setlength{\itemsep}{0.1ex}
        \item[$+$] $f$ pode ser caixa-preta
        \item[$+$] serve para $x$ discreto
        \item[$-$] variância alta
      \end{itemize}
    \end{defbox}
  \end{column}
  \begin{column}{0.5\textwidth}
    \begin{defbox}{Reparametrização}
      \vspace{-0.8em}
      \[ x = h_{\thv}(\epsilon),\; \gth \E[f] = \E\bigl[\nabla f \cdot \gth h_{\thv}\bigr] \]
      \vspace{-1.6em}
      \begin{itemize}\setlength{\itemsep}{0.1ex}
        \item[$+$] variância bem menor
        \item[$-$] exige $f$ diferenciável
        \item[$-$] só para $x$ contínuo
      \end{itemize}
    \end{defbox}
  \end{column}
  \end{columns}

  \begin{ideiabox}
    Em RL a ``função'' é o \emph{mundo}: o placar de um jogo, a preferência de
    um avaliador humano. Não há gradiente para pegar --- por isso o estimador
    por escore, caro em variância e barato em suposições, é o que sobrevive.
  \end{ideiabox}
\end{frame}

%% ------------------------------------------------------------
\begin{frame}{Teorema do gradiente de política}

  \begin{teobox}{Teorema \textnormal{(Sutton et al., 2000)}}
    Para toda política diferenciável $\polth(a\mid s)$ e para qualquer um dos
    objetivos usuais --- $J_1$ (recompensa episódica), $J_{\text{avR}}$
    (recompensa média por passo) ou $\tfrac{1}{1-\gamma}J_{\text{avV}}$ (valor
    médio) --- o gradiente da política é
    \[
      \gth J(\thv) = \E_{\polth}\bigl[\, \gscore\; Q^{\polth}(s,a) \,\bigr]
    \]
  \end{teobox}

  \begin{itemize}\setlength{\itemsep}{0.35ex}
    \item Generaliza o que derivamos: no lugar do retorno da trajetória
      inteira, aparece o valor da ação $Q^{\polth}(s,a)$.
    \item \textbf{O que o teorema evita:} a distribuição de estados
      $d^{\polth}$ também depende de $\thv$, mas o gradiente \emph{não} contém
      $\gth d^{\polth}$ --- esse é o resultado difícil (S\&B, cap.\ 13.2).
    \item Todo o resto da aula consiste em escolher \emph{o que colocar no
      lugar de $Q^{\polth}(s,a)$}.
  \end{itemize}
\end{frame}

%% ------------------------------------------------------------
\begin{frame}{Teste sua compreensão}

  \begin{quizbox}{}
    Sobre o gradiente por razão de verossimilhança, quais afirmações são
    verdadeiras?
    \begin{enumerate}\setlength{\itemsep}{0.1ex}
      \item[(a)] exige funções de recompensa diferenciáveis;
      \item[(b)] só pode ser usado com processos de decisão de Markov;
      \item[(c)] é útil principalmente em tarefas de horizonte infinito.
    \end{enumerate}
    \paradiscussao
  \end{quizbox}

  \pause
  \vspace{-0.3em}
  \begin{respbox}
    \textbf{Nenhuma das três.} (a) Nunca derivamos $\rec$, só $\log\polth$.
    (b) A derivação só usou que $\Ptraj$ fatora: vale para POMDPs, bandidos e
    geração de texto. (c) Ao contrário: supusemos episódios finitos.
  \end{respbox}
\end{frame}

%% ------------------------------------------------------------
\section{Reduzindo a variância}

%% ------------------------------------------------------------
\begin{frame}{O problema: não enviesado, mas ruidoso demais}

  \[
    \gth V(\thv) \approx \frac{1}{m}\sum_{i=1}^{m} R(\traj^{(i)})
      \sum_{t=0}^{T-1} \gth \log \polth\bigl(a_t^{(i)}\bigm| s_t^{(i)}\bigr)
  \]

  \begin{alertabox}
    Não enviesado --- e \textbf{de variância altíssima}. O sinal que atribuímos
    a \emph{cada} ação é o retorno da trajetória \emph{inteira}: uma ação
    excelente no passo 3 recebe a culpa por um desastre no passo 300.
  \end{alertabox}

  Três correções, nesta ordem, e cada uma é uma ideia independente:
  \begin{enumerate}\setlength{\itemsep}{0.35ex}
    \item \textbf{Estrutura temporal} --- uma ação não pode influenciar o
      passado.
    \item \textbf{Linha de base} --- comparar o retorno com o esperado, não
      com zero.
    \item \textbf{Alternativas ao retorno de Monte Carlo} --- trocar variância
      por um pouco de viés (ator--crítico).
  \end{enumerate}
\end{frame}

%% ------------------------------------------------------------
\begin{frame}{Correção 1: usar a estrutura temporal}

  Partindo de
  \[
    \gth \E_{\traj}[R] = \E_{\traj}\Bigl[\Bigl(\sum_{t=0}^{T-1} r_t\Bigr)
      \Bigl(\sum_{t=0}^{T-1} \gth \log \polth(a_t\mid s_t)\Bigr)\Bigr]
  \]

  \onslide<2->{
  O mesmo argumento, aplicado a \emph{uma única} recompensa $r_{t'}$, dá
  \[
    \gth \E[r_{t'}] = \E\Bigl[\, r_{t'} \sum_{t=0}^{t'} \gth \log \polth(a_t\mid s_t)\Bigr]
  \]
  --- só as ações \emph{até} $t'$ podem ter influenciado $r_{t'}$.}

  \onslide<3->{
  Somando em $t'$ e trocando a ordem dos somatórios:
  \[
    \gth V(\thv)
      = \E\Bigl[\sum_{t'=0}^{T-1} r_{t'} \sum_{t=0}^{t'} \gth \log \polth(a_t\mid s_t)\Bigr]
      = \E\Bigl[\sum_{t=0}^{T-1} \gth \log \polth(a_t\mid s_t)
        \mdestaque{rlLaranja}{\sum_{t'=t}^{T-1} r_{t'}}\Bigr]
  \]}
\end{frame}

%% ------------------------------------------------------------
\begin{frame}{Correção 1: o resultado}

  \begin{teobox}{Cada ação é creditada apenas pelo que veio depois}
    Como $\sum_{t'=t}^{T-1} r^{(i)}_{t'}$ é exatamente o retorno $G^{(i)}_t$,
    \[
      \gth V(\thv) \approx \frac{1}{m}\sum_{i=1}^{m}\sum_{t=0}^{T-1}
        \gth \log \polth\bigl(a^{(i)}_t\bigm| s^{(i)}_t\bigr)\, G^{(i)}_t
    \]
  \end{teobox}

  \begin{itemize}\setlength{\itemsep}{0.4ex}
    \item Removemos os termos $r_{t'}$ com $t' < t$: em esperança eles valiam
      \textbf{zero} (a ação não influencia o passado), mas contribuíam com
      variância pura.
    \item Continua não enviesado, com variância estritamente menor.
  \end{itemize}

  \begin{ideiabox}
    Princípio geral da aula: \emph{remova do estimador tudo que tem esperança
    zero}. Não muda o alvo e reduz ruído --- a linha de base, a seguir, é a
    mesma ideia.
  \end{ideiabox}
\end{frame}

%% ------------------------------------------------------------
\begin{frame}{REINFORCE: gradiente de política por Monte Carlo}

  \begin{algbox}{REINFORCE \textnormal{(Williams, 1992)}}
    {\small
    \begin{itemize}\setlength{\itemsep}{0.15ex}
      \item Inicialize $\thv$ arbitrariamente
      \item \textbf{para} cada episódio
        $\{s_0,a_0,r_0,\dots,s_{T-1},a_{T-1},r_{T-1}\} \sim \polth$ \textbf{faça}
        \begin{itemize}\setlength{\itemsep}{0.1ex}
          \item \textbf{para} $t = 0$ até $T-1$ \textbf{faça}
            \[ \thv \leftarrow \thv + \alpha\, \gth \log \polth(a_t\mid s_t)\, G_t \]
          \item \textbf{fim para}
        \end{itemize}
      \item \textbf{fim para}; devolva $\thv$
    \end{itemize}}
  \end{algbox}

  \begin{itemize}\setlength{\itemsep}{0.3ex}
    \item Combina razão de verossimilhança $+$ estrutura temporal.
    \item É \textbf{on-policy} e exige o episódio completo (precisa de $G_t$).
    \item Publicado em 1992 --- é o ancestral direto de tudo que vem depois.
  \end{itemize}
\end{frame}

%% ------------------------------------------------------------
\begin{frame}{Por que ainda não funciona bem}{Um exemplo de três linhas}

  Bandido de duas ações, política softmax, recompensas
  $\rec(a_1) = 100$ e $\rec(a_2) = 101$.

  \begin{itemize}\setlength{\itemsep}{0.35ex}
    \item Ambos os retornos são grandes e positivos $\Rightarrow$ o estimador
      \textbf{aumenta a log-probabilidade das duas ações} em toda amostra.
    \item A diferença entre elas --- o único sinal útil --- é $1$, sepultada
      num sinal de magnitude $100$.
    \item Agora some $+1000$ a todas as recompensas. O gradiente
      \emph{verdadeiro} não muda (o escore tem média zero), mas a
      \textbf{variância do estimador explode}.
  \end{itemize}

  \begin{ideiabox}
    A escala absoluta da recompensa não deveria importar --- e não importa em
    esperança. Ela só afeta o \emph{ruído}. Consertar isso é exatamente o
    papel da linha de base.
  \end{ideiabox}
\end{frame}

%% ------------------------------------------------------------
\begin{frame}{Correção 2: introduzir uma linha de base}

  \begin{defbox}{Linha de base $b(s)$}
    \vspace{-0.9em}
    \[
      \gth \E_{\traj}[R] = \E_{\traj}\Bigl[\sum_{t=0}^{T-1}
        \gth \log \polth(a_t\mid s_t)\bigl(G_t - b(s_t)\bigr)\Bigr]
    \]
    \vspace{-0.9em}

    Para \textbf{qualquer} $b$ que dependa apenas do estado (não da ação), o
    estimador permanece não enviesado.
  \end{defbox}

  \begin{itemize}\setlength{\itemsep}{0.35ex}
    \item Escolha quase ótima: o retorno esperado,
      $b(s_t) \approx \E[r_t + \dots + r_{T-1}] = \Vpi(s_t)$.
    \item \textbf{Leitura:} aumente a log-probabilidade de $a_t$ na proporção
      de quanto o retorno foi \emph{melhor que o esperado} naquele estado.
  \end{itemize}

  \begin{ideiabox}
    De ``esta ação rendeu 101'' para ``esta ação rendeu $0{,}5$ acima da média
    deste estado''. O segundo sinal é comparável entre estados; o primeiro, não.
  \end{ideiabox}
\end{frame}

%% ------------------------------------------------------------
\begin{frame}{A linha de base não introduz viés --- demonstração}

  {\small
  \begin{align*}
    \E_{\traj}\bigl[\gth \log \polth(a_t\mid s_t)\, b(s_t)\bigr]
      &= \E_{s_{0:t},a_{0:t-1}}\Bigl[\E_{a_t}\bigl[\gth \log \polth(a_t\mid s_t)\, b(s_t)\bigr]\Bigr]
        &&\text{\scriptsize (quebrar a esperança)}\\
      &= \E_{s_{0:t},a_{0:t-1}}\Bigl[b(s_t)\,\E_{a_t}\bigl[\gth \log \polth(a_t\mid s_t)\bigr]\Bigr]
        &&\text{\scriptsize ($b$ não depende de $a_t$)}\\
      &= \E\Bigl[b(s_t) \sum_{a} \polth(a\mid s_t)
         \frac{\gth \polth(a\mid s_t)}{\polth(a\mid s_t)}\Bigr]
        &&\text{\scriptsize (razão de verossimilhança)}\\
      &= \E\Bigl[b(s_t)\, \gth \sum_{a} \polth(a\mid s_t)\Bigr]
      \;=\; \E\bigl[b(s_t)\, \gth 1\bigr] \;=\; 0
  \end{align*}}
  \vspace{-0.6em}

  \begin{teobox}{Consequência}
    Subtrair $b(s_t)$ altera a variância mas \textbf{não} a esperança: há
    liberdade total para escolher $b$ minimizando o ruído. Atenção ao segundo
    passo --- ele exige que $b$ dependa só de $s_t$.
  \end{teobox}
\end{frame}

%% ------------------------------------------------------------
\begin{frame}{O algoritmo ``baunilha'' de gradiente de política}

  \begin{algbox}{Gradiente de política com linha de base}
    {\footnotesize
    \begin{itemize}\setlength{\itemsep}{0.15ex}
      \item Inicialize os parâmetros da política $\thv$ e a linha de base $b$
      \item \textbf{para} iteração $=1,2,\dots$ \textbf{faça}
        \begin{itemize}\setlength{\itemsep}{0.1ex}
          \item Colete trajetórias executando a política atual
          \item Em cada passo $t$ de cada trajetória $\traj^i$, calcule o
            retorno $G^i_t = \sum_{t'=t}^{T-1} r^i_{t'}$ e a vantagem
            $\Ahat^i_t = G^i_t - b(s_t)$
          \item \textbf{Reajuste a linha de base} minimizando
            $\sum_i \sum_t \norm{b(s_t) - G^i_t}^2$
          \item \textbf{Atualize a política} com
            $\hat{g} = \sum_{i,t} \gth \log \polth(a_t\mid s_t)\, \Ahat^i_t$
            (via SGD ou Adam)
        \end{itemize}
      \item \textbf{fim para}
    \end{itemize}}
  \end{algbox}

  \vspace{-0.6em}
  \begin{ideiabox}
    Dois problemas supervisionados alternando: regressão (ajustar $b$) e
    subida de gradiente ponderada (ajustar $\polth$).
  \end{ideiabox}
\end{frame}

%% ------------------------------------------------------------
\begin{frame}{Que linha de base escolher?}

  \begin{itemize}\setlength{\itemsep}{0.45ex}
    \item \textbf{Constante:} $b = $ média dos retornos do lote. Trivial,
      já ajuda muito.
    \item \textbf{Dependente do estado:} $b(s) = \Vhat^{\pol}(s;\wv)$, uma
      rede treinada por regressão. É a escolha padrão.
    \item \textbf{Ótima em variância:} a linha de base que minimiza
      $\Var(\hat g)$ é uma média dos retornos \emph{ponderada pelo escore ao
      quadrado},
      \[
        b^{*} = \frac{\E\bigl[(\gth\log\polth)^2 G_t\bigr]}
                     {\E\bigl[(\gth\log\polth)^2\bigr]}
      \]
      Raramente vale o custo: $\Vhat^{\pol}$ chega perto e é mais simples.
  \end{itemize}

  \begin{alertabox}
    Na prática faz-se também \textbf{normalização de vantagens} por lote:
    $\Ahat \leftarrow (\Ahat - \bar{\Ahat})/(\sigma_{\Ahat} + \varepsilon)$ ---
    viés pequeno, e o passo $\alpha$ deixa de depender da escala da recompensa.
  \end{alertabox}
\end{frame}

%% ------------------------------------------------------------
\begin{frame}{Teste sua compreensão}

  \begin{quizbox}{}
    Um aluno propõe usar como linha de base
    $b(s_t,a_t) = \Qhat^{\pol}(s_t,a_t)$ --- afinal, dizem os slides, a
    escolha quase ótima é ``o retorno esperado'', e $\Qhat$ estima o retorno
    esperado com mais precisão que $\Vhat$.

    \smallskip
    O estimador resultante ainda é não enviesado? O que ele calcularia?
    \paradiscussao
  \end{quizbox}

  \pause

  \begin{respbox}
    \textbf{Não.} A demonstração de não-viés exige que $b$ dependa apenas de
    $s_t$; com $b$ dependendo de $a_t$, não se pode retirá-la da esperança
    sobre ações. Pior: o termo subtraído seria exatamente
    $\E[\gth\log\pol \cdot Q^{\pol}]$, ou seja, o gradiente inteiro --- em
    esperança o estimador daria \textbf{zero} e a política nunca aprenderia.
  \end{respbox}
\end{frame}

%% ------------------------------------------------------------
\section{Alternativas ao retorno de Monte Carlo}

%% ------------------------------------------------------------
\begin{frame}{Onde estamos e o que falta}

  \[
    \gth \E[R] \approx \frac{1}{m}\sum_{i=1}^{m}\sum_{t=0}^{T-1}
      \gth \log \polth(a_t\mid s_t)\bigl(G^{(i)}_t - b(s_t)\bigr)
  \]

  \begin{itemize}\setlength{\itemsep}{0.15ex}
    \item[\checkmark] Estrutura temporal e linha de base --- ambas no
      somatório acima.
    \item[$\square$] \textbf{Alternativas a $G_t$} como estimativa do retorno
      esperado.
  \end{itemize}

  \begin{alertabox}
    $G^i_t$ estima o valor a partir de \emph{uma única} execução: não
    enviesada, mas carregando o ruído de todas as transições e escolhas de
    ação até o fim do episódio.
  \end{alertabox}

  \begin{ideiabox}
    O remédio é o da Aula 4, quando comparamos MC e TD: introduzir
    \textbf{viés} via \emph{bootstrapping} para comprar uma redução grande de
    \textbf{variância}.
  \end{ideiabox}
\end{frame}

%% ------------------------------------------------------------
\begin{frame}{Métodos ator--crítico}

  \begin{defbox}{Ator e crítico}
    O \textbf{ator} é a política $\polth$ --- decide o que fazer. O
    \textbf{crítico} é a função valor $\Vhat(s;\wv)$ ou $\Qhat(s,a;\wv)$ ---
    avalia o que o ator fez. Ambos são aprendidos e atualizados.
  \end{defbox}

  \begin{center}
  \scalebox{0.86}{\begin{tikzpicture}[font=\scriptsize, node distance=8mm]
    \node[draw=rlLaranja, very thick, fill=rlLaranja!10, rounded corners=3pt,
          minimum width=26mm, minimum height=9mm] (ator) {\textbf{Ator} $\polth(a\mid s)$};
    \node[draw=rlTeal, very thick, fill=rlTeal!10, rounded corners=3pt,
          minimum width=26mm, minimum height=9mm, below=13mm of ator] (crit)
          {\textbf{Crítico} $\Vhat(s;\wv)$};
    \node[draw=rlAzul, very thick, fill=rlAzul!8, rounded corners=3pt,
          minimum width=26mm, minimum height=9mm, right=34mm of ator] (amb) {\textbf{Ambiente}};
    \draw[trans destac] (ator) -- node[above] {ação $a_t$} (amb);
    \draw[trans] (amb.south) |- node[pos=0.75, above] {$r_t,\, s_{t+1}$} (crit.east);
    \draw[trans] (crit) -- node[right] {vantagem $\Ahat_t$} (ator);
  \end{tikzpicture}}
  \end{center}

  \begin{itemize}\setlength{\itemsep}{0.25ex}
    \item O crítico troca o retorno bruto por uma estimativa
      \emph{bootstrapped}: menos ruidosa, ligeiramente enviesada. A3C (Mnih et
      al., 2016) popularizou a arquitetura; A2C, PPO e SAC a herdam.
  \end{itemize}
\end{frame}

%% ------------------------------------------------------------
\begin{frame}{A função vantagem}

  Partindo de
  $\gth \E_\traj[R] \approx \E\bigl[\sum_t \gth\log\polth(a_t\mid s_t)\,
   \bigl(Q(s_t,a_t;\wv) - b(s_t)\bigr)\bigr]$
  e tomando $b(s) = \Vhat^{\pol}(s)$:

  \begin{defbox}{Vantagem}
    \[
      \Advpi(s,a) = Q^{\pol}(s,a) - \Vpi(s)
      \qquad\Longrightarrow\qquad
      \gth \E_\traj[R] \approx \E\Bigl[\sum_{t=0}^{T-1}
        \gth \log \polth(a_t\mid s_t)\, \Ahat^{\pol}(s_t,a_t)\Bigr]
    \]
    ``Quanto esta ação é melhor do que a ação média neste estado.''
  \end{defbox}

  \begin{itemize}\setlength{\itemsep}{0.3ex}
    \item Por construção, $\E_{a\sim\pol}[\Advpi(s,a)] = 0$: a vantagem é
      centrada.
    \item Positivo $\Rightarrow$ torne a ação mais provável; negativo
      $\Rightarrow$ menos. Sem esse centramento, ``todas as ações ficam mais
      prováveis'' --- o que não faz sentido para uma distribuição.
  \end{itemize}
\end{frame}

%% ------------------------------------------------------------
\begin{frame}{Estimadores de $n$ passos: um dial entre TD e MC}

  O crítico pode usar \textbf{qualquer mistura} entre alvo TD e alvo MC:

  \[
    \hat{R}^{(1)}_t = r_t + \gamma \Vhat(s_{t+1}),
    \quad
    \hat{R}^{(2)}_t = r_t + \gamma r_{t+1} + \gamma^2 \Vhat(s_{t+2}),
    \quad \dots, \quad
    \hat{R}^{(\infty)}_t = \sum_{k\ge 0}\gamma^k r_{t+k}
  \]

  Subtraindo a linha de base $\Vhat(s_t)$, obtêm-se estimadores de vantagem:
  \[
    \Ahat^{(1)}_t = r_t + \gamma \Vhat(s_{t+1}) - \Vhat(s_t),
    \qquad
    \Ahat^{(\infty)}_t = \sum_{k\ge 0}\gamma^k r_{t+k} - \Vhat(s_t)
  \]

  \begin{ideiabox}
    $\Ahat^{(1)}$ é exatamente o \textbf{erro TD} $\delta_t$ da Aula 4, agora
    reaparecendo como sinal de aprendizado da \emph{política}. O mesmo objeto,
    dois papéis.
  \end{ideiabox}
\end{frame}

%% ------------------------------------------------------------
\begin{frame}{O compromisso viés--variância}

  \begin{center}
  \scalebox{0.95}{\begin{tikzpicture}[font=\scriptsize]
    \draw[rlCinza!40, line width=7pt] (0,0) -- (10,0);
    \foreach \x/\lab in {0/{$\Ahat^{(1)}$}, 2.5/{$\Ahat^{(2)}$}, 5/{$\Ahat^{(4)}$},
                         7.5/{$\Ahat^{(8)}$}, 10/{$\Ahat^{(\infty)}$}}{
      \fill[rlAzul] (\x,0) circle (2.6pt);
      \node[below=2mm, rlAzul] at (\x,0) {\lab};}
    \node[above=6mm, rlTeal, align=center] at (0,0) {\textbf{TD puro}\\ 1 passo};
    \node[above=6mm, rlLaranja, align=center] at (10,0) {\textbf{MC puro}\\ episódio inteiro};
    \draw[rlVermelho, very thick, -{Stealth[length=2.2mm]}] (0.3,1.55) -- (9.7,1.55);
    \node[rlVermelho, above] at (5,1.55) {variância cresce};
    \draw[rlAmbar!80!black, very thick, -{Stealth[length=2.2mm]}] (9.7,-1.5) -- (0.3,-1.5);
    \node[rlAmbar!80!black, below] at (5,-1.5) {viés cresce};
  \end{tikzpicture}}
  \end{center}

  \begin{itemize}\setlength{\itemsep}{0.3ex}
    \item $n$ pequeno: alvo depende fortemente de $\Vhat$, que está errado no
      início $\Rightarrow$ \textbf{viés}.
    \item $n$ grande: alvo acumula o ruído de muitas transições $\Rightarrow$
      \textbf{variância}.
    \item \textbf{GAE} (Schulman et al., 2016) toma uma média exponencial de
      todos os $n$ com peso $\lambda$:
      $\Ahat^{\text{GAE}}_t = \sum_{k\ge 0}(\gamma\lambda)^k \delta_{t+k}$.
      É o padrão em PPO --- e é exatamente o $\lambda$ de TD($\lambda$).
  \end{itemize}
\end{frame}

%% ------------------------------------------------------------
\begin{frame}{Teste sua compreensão}

  \begin{quizbox}{}
    Considerando
    \[
      \Ahat^{(1)}_t = r_t + \gamma \Vhat(s_{t+1}) - \Vhat(s_t),
      \qquad
      \Ahat^{(\infty)}_t = r_t + \gamma r_{t+1} + \gamma^2 r_{t+2} + \cdots - \Vhat(s_t)
    \]
    Classifique cada um quanto a viés e variância.
    \paradiscussao
  \end{quizbox}

  \pause

  \begin{respbox}
    $\Ahat^{(1)}$ tem \textbf{variância baixa e viés alto}: depende de uma
    única transição, mas herda todo o erro de $\Vhat(s_{t+1})$.

    \smallskip
    $\Ahat^{(\infty)}$ tem \textbf{variância alta e viés baixo}: usa
    recompensas reais até o fim (só $\Vhat(s_t)$, a linha de base, é
    aproximada --- e ela não introduz viés), ao custo de acumular o ruído de
    todo o episódio.
  \end{respbox}
\end{frame}

%% ------------------------------------------------------------
\begin{frame}{O algoritmo baunilha, versão final}

  \begin{algbox}{Gradiente de política com vantagem}
    {\footnotesize
    \begin{itemize}\setlength{\itemsep}{0.15ex}
      \item Inicialize os parâmetros da política $\thv$ e do crítico $\wv$
      \item \textbf{para} iteração $=1,2,\dots$ \textbf{faça}
        \begin{itemize}\setlength{\itemsep}{0.1ex}
          \item Colete trajetórias executando $\polth$
          \item Em cada passo, calcule uma estimativa de vantagem
            $\Ahat^{(n)}_{it}$ (escolha de $n$, ou GAE$(\lambda)$)
          \item Ajuste o crítico: $\wv \leftarrow \argmin_{\wv}
            \sum_{i,t}\bigl(\Vhat(s_{it};\wv) - \hat{R}_{it}\bigr)^2$
          \item Atualize a política com
            $\hat{g} = \sum_{i,t}\gth \log \polth(a_t\mid s_t)\,\Ahat^{(n)}_{it}$
        \end{itemize}
      \item \textbf{fim para}
    \end{itemize}}
  \end{algbox}

  \begin{ideiabox}
    Toda a família --- REINFORCE, A2C, TRPO, PPO, GRPO --- é este esqueleto
    com duas escolhas: \emph{como estimar a vantagem} e \emph{como limitar o
    passo}. A próxima aula trata da segunda.
  \end{ideiabox}
\end{frame}

%% ------------------------------------------------------------
\section{Balanço e próximos passos}

%% ------------------------------------------------------------
\begin{frame}{Vantagens e desvantagens do RL baseado em política}

  \begin{columns}[T]
  \begin{column}{0.5\textwidth}
    \begin{defbox}{Vantagens}
      \begin{itemize}\setlength{\itemsep}{0.2ex}
        \item Convergência suave, sem oscilação por troca de $\argmax$.
        \item Eficaz em ações \textbf{contínuas} ou de alta dimensão.
        \item Aprende políticas \textbf{estocásticas}.
      \end{itemize}
    \end{defbox}
  \end{column}
  \begin{column}{0.5\textwidth}
    \begin{alertabox}
      \begin{itemize}\setlength{\itemsep}{0.2ex}
        \item Converge a um \textbf{ótimo local}, não global.
        \item Avaliar a política é ineficiente e de \textbf{alta variância}.
        \item É \textbf{on-policy}: cada lote serve a uma atualização e é
          descartado.
      \end{itemize}
    \end{alertabox}
  \end{column}
  \end{columns}

  \begin{ideiabox}
    Escolha prática: ações discretas e simulador caro $\Rightarrow$ DQN. Ações
    contínuas ou política necessariamente estocástica $\Rightarrow$ gradiente
    de política. Os dois $\Rightarrow$ ator--crítico off-policy (SAC, DDPG).
  \end{ideiabox}
\end{frame}

%% ------------------------------------------------------------
\begin{frame}{Armadilhas práticas}

  \begin{itemize}\setlength{\itemsep}{0.45ex}
    \item \textbf{Passo $\alpha$ é criticamente sensível.} Um passo grande
      demais destrói a política; como os dados seguintes são coletados
      \emph{por ela}, não há como se recuperar. Este é o problema que motiva
      TRPO/PPO na próxima aula.
    \item \textbf{Colapso de entropia.} A política fica determinística cedo,
      a exploração morre, e o gradiente vai a zero. Remédio usual: um bônus de
      entropia $+\beta\,\mathcal{H}[\polth(\cdot\mid s)]$ no objetivo.
    \item \textbf{Normalize as vantagens} por lote; normalize também as
      observações.
    \item \textbf{Lotes grandes.} Gradiente de política precisa de muito mais
      amostras por atualização do que aprendizado supervisionado --- a
      variância cai com $1/\sqrt{m}$.
    \item \textbf{Recompensas esparsas} continuam sendo o pior caso: se
      nenhuma trajetória do lote obtém recompensa, $\hat g \approx 0$ e nada
      acontece.
  \end{itemize}
\end{frame}

%% ------------------------------------------------------------
\begin{frame}{O que você deve levar desta aula}

  \begin{defbox}{Capacidades esperadas}
    \begin{itemize}\setlength{\itemsep}{0.25ex}
      \item Explicar por que uma política estocástica pode ser melhor sob
        observabilidade parcial.
      \item \textbf{Derivar} o gradiente por razão de verossimilhança e dizer
        por que a dinâmica desaparece.
      \item Calcular o escore de políticas softmax e gaussiana; implementar
        REINFORCE com e sem linha de base.
      \item Demonstrar que $b(s)$ não introduz viés --- e o que quebra se ela
        depender da ação.
      \item Posicionar $\Ahat^{(n)}$ no compromisso viés--variância.
    \end{itemize}
  \end{defbox}

  \vspace{-0.9em}
  \begin{ideiabox}
    A moral: gradiente de política é ``aumente a probabilidade do que deu
    certo''; o resto foi definir \emph{o que conta como ter dado certo}.
  \end{ideiabox}
\end{frame}

%% ------------------------------------------------------------
\begin{frame}{Próxima aula}

  \begin{itemize}\setlength{\itemsep}{0.45ex}
    \item O estimador está pronto; o \emph{passo} não. Quão longe podemos
      andar em $\thv$ sem destruir a política que gera nossos dados?
    \item \textbf{Gradiente natural}: medir a distância entre políticas no
      espaço de \emph{distribuições} (divergência KL), não no espaço de
      parâmetros.
    \item \textbf{TRPO} e \textbf{PPO}: garantir melhoria monotônica
      aproximada limitando a KL --- e a versão barata, com razão de
      probabilidades recortada, que treina modelos de linguagem hoje.
  \end{itemize}

  \begin{center}
  \scalebox{0.9}{\begin{tikzpicture}[font=\scriptsize, node distance=7mm]
    \node[draw=rlCinza, thick, fill=rlFundo, rounded corners=2pt,
          minimum width=26mm, minimum height=8mm] (a) {REINFORCE};
    \node[draw=rlAzul, thick, fill=rlAzul!8, rounded corners=2pt,
          minimum width=26mm, minimum height=8mm, right=10mm of a] (b) {$+$ linha de base};
    \node[draw=rlTeal, thick, fill=rlTeal!10, rounded corners=2pt,
          minimum width=26mm, minimum height=8mm, right=10mm of b] (c) {ator--crítico};
    \node[draw=rlLaranja, very thick, fill=rlLaranja!12, rounded corners=2pt,
          minimum width=26mm, minimum height=8mm, right=10mm of c] (d) {TRPO / PPO};
    \draw[trans] (a) -- (b); \draw[trans] (b) -- (c); \draw[trans destac] (c) -- (d);
  \end{tikzpicture}}
  \end{center}

  \begin{ideiabox}
    Leitura: Sutton \& Barto (2018), cap.\ 13. Para o lado prático, o artigo
    de GAE (Schulman et al., 2016) é curto e vale cada página.
  \end{ideiabox}
\end{frame}

\end{document}
